手机上的 GUI 智能体在执行长任务时，常常要输入一段十几屏之前才见过的文字：它在菜谱页看到"香草精"，九步之后要把这三个字打进购物清单，而那一屏早已不在眼前。难点不在于历史存不存得下，而在于智能体得知道该回头看哪一屏。

要训练模型做这个判断，通常得有人标注"哪一屏是对的"。我们发现不必。这类题目本来就靠一个自动检查挖出来：输入的文字出现在某个更早屏的 OCR 文本里，而不在当前屏上。把这同一个检查施加到每一屏，就直接得到了每屏的对错标签，一个人也不用看。我们把用这种标签训练回看策略的做法称为\textbf{验证器监督的回看}。

我们用这份免费标签训练了两种模型：一种一次读完整段历史、直接说出屏号，另一种每次只看一屏、给出"这屏是不是"的概率。两者都大幅变好。在 188 道题上，前者找对屏的比例从 54.8\% 升到 65.4\%，后者从 50.0\% 升到 72.3\%；在一批事后才挖、只测一次的新轨迹上，它们分别达到 76.3\% 与 79.6\%，而不训练只给 prompt 的做法是 59.1\%。最终输入对文字的比例也从 23.7\% 升到约 50\%。

值得注意的是，两种模型训练后谁更好，统计上分不出来（$p=0.086$）。换句话说，\textbf{真正起作用的是这份免费标签，而不是我们一开始押注的逐屏架构}。逐屏的好处只在一个具体区间成立：当正确的那一屏在 8 到 12 屏之前时，它 78.2\% 对 56.4\%（\pval{0.003}）。我们也如实报告两种模型共同失效的地方——距离超过 13 屏时，它们都不如不训练的做法——以及任何人研究这个问题都必须先排掉的两个评测陷阱。
